% GENERATED FILE. Edit canonical lessons, then run npm run generate:books.
% canonical-source-hash: fnv1a64:8439e860986cfa72
% canonical-lessons: SA-C327-holika, SA-C327-kunkumam, SA-C327-sincati, SA-C327-vikirati, SA-C327-utksipati

\chapter{A Festival of Colours}
\label{ch:sa-a-festival-of-colours}
\hlchaptermodality{327}

\begin{chapteropening}
\textbf{By the end of this chapter:} \emph{I can tell how a spring festival was played.}

Complete the last lesson of chapter 327: I can tell how a spring festival was played.
\end{chapteropening}

\section[\texorpdfstring{\sk{होलिका} (holikā)}{होलिका (holikā)}]{\sk{होलिका} (holikā) — the festival of colours; Holi}

\label{lesson:SA-C327-holika}



\begin{quote}
Before the new one: say the Sanskrit for festive excitement, then the Sanskrit for splendour.
\end{quote}

\subsection*{You'll want to know: \sk{होलिका}}
\textbf{\sk{होलिका}} — \emph{holikā} — \textquotedblleft{}the festival of colours; Holi\textquotedblright{}.

\begin{quote}
\sk{वसन्ते} \sk{होलिका} \sk{भवति।}
\end{quote}

\begin{quote}
\emph{vasante holikā bhavati.}
\end{quote}

\begin{quote}
Holi comes in spring.
\end{quote}

A feminine noun.

\begin{cousinweb}
\textbf{\sk{होलिका}} (\emph{holikā}) MW gives \emph{holikā} as the Holi festival, beside \emph{holākā}, the spring festival near the equinox, which it takes perhaps from a cry or shout in singing, so its origin is uncertain (MW).
\end{cousinweb}

\begin{grammarlens}[title={What you've built}]
One more word for telling what happened.
\end{grammarlens}

\subsection*{Your turn}
\begin{itemize}
\raggedright
  \item \emph{Say it:} \emph{holikā}
  \item \emph{Say it:} \emph{holikā}, once more
  \item \emph{Say it:} \emph{vasante holikā bhavati.} — Holi comes in spring.
  \item \emph{Recall:} say the Sanskrit for festive excitement, then the Sanskrit for splendour, then say \emph{holikā} again
\end{itemize}

\begin{tcolorbox}[breakable,colback=teal!4,colframe=teal!35!black,title={Before you move on}]
What does \sk{होलिका} mean? (\textquotedblleft{}The festival of colours; Holi\textquotedblright{}.) Say it once more.
\end{tcolorbox}
\section[\texorpdfstring{\sk{कुंकुमम्} (kuṅkumam)}{कुंकुमम् (kuṅkumam)}]{\sk{कुंकुमम्} (kuṅkumam) — saffron}

\label{lesson:SA-C327-kunkumam}



\begin{quote}
Before the new one: say the Sanskrit for splendour, then the Sanskrit for the festival of colours.
\end{quote}

\subsection*{You'll want to know: \sk{कुंकुमम्}}
\textbf{\sk{कुंकुमम्}} — \emph{kuṅkumam} — \textquotedblleft{}saffron\textquotedblright{}.

\begin{quote}
\sk{पायसे} \sk{कुंकुमम्} \sk{अस्ति।}
\end{quote}

\begin{quote}
\emph{pāyase kuṅkumam asti.}
\end{quote}

\begin{quote}
There is saffron in the rice pudding.
\end{quote}

A neuter noun.

\begin{cousinweb}
\textbf{\sk{कुंकुमम्}} (\emph{kuṅkumam}) MW gives \emph{kuṅkuma} as saffron, the plant and the pollen of its flowers, without a derivation, so its origin is uncertain (MW).
\end{cousinweb}

\begin{grammarlens}[title={What you've built}]
One more word for telling what happened.
\end{grammarlens}

\subsection*{Your turn}
\begin{itemize}
\raggedright
  \item \emph{Say it:} \emph{kuṅkumam}
  \item \emph{Say it:} \emph{kuṅkumam}, once more
  \item \emph{Say it:} \emph{pāyase kuṅkumam asti.} — There is saffron in the rice pudding.
  \item \emph{Recall:} say the Sanskrit for splendour, then the Sanskrit for the festival of colours, then say \emph{kuṅkumam} again
\end{itemize}

\begin{tcolorbox}[breakable,colback=teal!4,colframe=teal!35!black,title={Before you move on}]
What does \sk{कुंकुमम्} mean? (\textquotedblleft{}Saffron\textquotedblright{}.) Say it once more.
\end{tcolorbox}
\section[\texorpdfstring{\sk{सिञ्चति} (siñcati)}{सिञ्चति (siñcati)}]{\sk{सिञ्चति} (siñcati) — sprinkles}

\label{lesson:SA-C327-sincati}



\begin{quote}
Before the new one: say the Sanskrit for the festival of colours, then the Sanskrit for saffron.
\end{quote}

\subsection*{You'll want to know: \sk{सिञ्चति}}
\textbf{\sk{सिञ्चति}} — \emph{siñcati} — \textquotedblleft{}sprinkles\textquotedblright{}.

\begin{quote}
\sk{बालाः} \sk{मित्रेषु} \sk{जलं} \sk{सिञ्चन्ति।}
\end{quote}

\begin{quote}
\emph{bālāḥ mitreṣu jalaṃ siñcanti.}
\end{quote}

\begin{quote}
The children sprinkle water on their friends.
\end{quote}

Past (he …): \sk{सः} \sk{मित्रे} \sk{जलं} \sk{सिक्तवान्}, \emph{saḥ mitre jalaṃ siktavān}.

\begin{cousinweb}
\textbf{\sk{सिञ्चति}} (\emph{siñcati}) comes from the root √sic, \textquotedblleft{}to pour out, sprinkle\textquotedblright{}, the root of \emph{abhiṣeka} too (MW).
\end{cousinweb}

\begin{grammarlens}[title={What you've built}]
One more word for telling what happened.
\end{grammarlens}

\subsection*{Your turn}
\begin{itemize}
\raggedright
  \item \emph{Say it:} \emph{siñcati}
  \item \emph{Say it:} \emph{siñcati}, once more
  \item \emph{Say it:} \emph{bālāḥ mitreṣu jalaṃ siñcanti.} — The children sprinkle water on their friends.
  \item \emph{Recall:} say the Sanskrit for the festival of colours, then the Sanskrit for saffron, then say \emph{siñcati} again
\end{itemize}

\begin{tcolorbox}[breakable,colback=teal!4,colframe=teal!35!black,title={Before you move on}]
What does \sk{सिञ्चति} mean? (\textquotedblleft{}Sprinkles\textquotedblright{}.) Say it once more.
\end{tcolorbox}
\section[\texorpdfstring{\sk{विकिरति} (vikirati)}{विकिरति (vikirati)}]{\sk{विकिरति} (vikirati) — scatters}

\label{lesson:SA-C327-vikirati}



\begin{quote}
Before the new one: say the Sanskrit for saffron, then the Sanskrit for sprinkles.
\end{quote}

\subsection*{You'll want to know: \sk{विकिरति}}
\textbf{\sk{विकिरति}} — \emph{vikirati} — \textquotedblleft{}scatters\textquotedblright{}.

\begin{quote}
\sk{जनाः} \sk{मार्गे} \sk{पुष्पाणि} \sk{विकिरन्ति।}
\end{quote}

\begin{quote}
\emph{janāḥ mārge puṣpāṇi vikiranti.}
\end{quote}

\begin{quote}
People scatter flowers on the road.
\end{quote}

Past (he …): \sk{सः} \sk{पुष्पाणि} \sk{विकीर्णवान्}, \emph{saḥ puṣpāṇi vikīrṇavān}.

\begin{cousinweb}
\textbf{\sk{विकिरति}} (\emph{vikirati}) is built from \emph{vi} and the root √kṝ, \textquotedblleft{}to scatter\textquotedblright{}; MW gives \emph{vi-kṝ} as to scatter, throw about, disperse (MW).
\end{cousinweb}

\begin{grammarlens}[title={What you've built}]
One more word for telling what happened.
\end{grammarlens}

\subsection*{Your turn}
\begin{itemize}
\raggedright
  \item \emph{Say it:} \emph{vikirati}
  \item \emph{Say it:} \emph{vikirati}, once more
  \item \emph{Say it:} \emph{janāḥ mārge puṣpāṇi vikiranti.} — People scatter flowers on the road.
  \item \emph{Recall:} say the Sanskrit for saffron, then the Sanskrit for sprinkles, then say \emph{vikirati} again
\end{itemize}

\begin{tcolorbox}[breakable,colback=teal!4,colframe=teal!35!black,title={Before you move on}]
What does \sk{विकिरति} mean? (\textquotedblleft{}Scatters\textquotedblright{}.) Say it once more.
\end{tcolorbox}
\section[\texorpdfstring{\sk{उत्क्षिपति} (utkṣipati)}{उत्क्षिपति (utkṣipati)}]{\sk{उत्क्षिपति} (utkṣipati) — tosses up}

\label{lesson:SA-C327-utksipati}



\begin{quote}
Before the new one: say the Sanskrit for sprinkles, then the Sanskrit for scatters.
\end{quote}

\subsection*{You'll want to know: \sk{उत्क्षिपति}}
\textbf{\sk{उत्क्षिपति}} — \emph{utkṣipati} — \textquotedblleft{}tosses up\textquotedblright{}.

\begin{quote}
\sk{बालः} \sk{कन्दुकम्} \sk{उपरि} \sk{उत्क्षिपति।}
\end{quote}

\begin{quote}
\emph{bālaḥ kandukam upari utkṣipati.}
\end{quote}

\begin{quote}
The boy tosses the ball up.
\end{quote}

Past (he …): \sk{सः} \sk{कन्दुकम्} \sk{उत्क्षिप्तवान्}, \emph{saḥ kandukam utkṣiptavān}.

\begin{cousinweb}
\textbf{\sk{उत्क्षिपति}} (\emph{utkṣipati}) is built from \emph{ud}, \textquotedblleft{}up\textquotedblright{}, and the root √kṣip, \textquotedblleft{}to throw\textquotedblright{}, as in \emph{kṣipati}; MW gives \emph{ut-kṣip} as to throw up, raise (MW).
\end{cousinweb}

\begin{grammarlens}[title={What you've built}]
One more word for telling what happened.
\end{grammarlens}

\subsection*{Your turn}
\begin{itemize}
\raggedright
  \item \emph{Say it:} \emph{utkṣipati}
  \item \emph{Say it:} \emph{utkṣipati}, once more
  \item \emph{Say it:} \emph{bālaḥ kandukam upari utkṣipati.} — The boy tosses the ball up.
  \item \emph{Recall:} say the Sanskrit for sprinkles, then the Sanskrit for scatters, then say \emph{utkṣipati} again
\end{itemize}

\begin{tcolorbox}[breakable,colback=teal!4,colframe=teal!35!black,title={Before you move on}]
What does \sk{उत्क्षिपति} mean? (\textquotedblleft{}Tosses up\textquotedblright{}.) Say it once more.
\end{tcolorbox}
